\section{Related Work}
\label{sec:related_work}

\subsection{USD and MDL in Robotics Simulation}

Universal Scene Description (USD) has emerged as the standard scene format for
robotics simulation, underpinning NVIDIA's Isaac
Lab~\cite{Mittal2025Isaac} and its predecessor
Orbit~\cite{Mittal2023Orbit}. These platforms combine USD's compositional scene
graph with MDL-based physically accurate materials rendered via RTX ray tracing.
Downstream pipelines such as SynTable~\cite{Ng2023SynTable} and
Synthetica~\cite{Singh2024Synthetica} generate millions of labeled images using
this rendering stack, demonstrating its importance for training perception
models. However, none of these works examine the consequences of simplifying
MDL to lighter-weight material representations, a gap our study directly
addresses.

\subsection{Sim-to-Real Visual Gap}

The disparity between simulated and real images has been studied extensively.
Domain randomization~\cite{Tobin2017Domain,Tremblay2018Training} improves
transfer by diversifying visual parameters such as textures and lighting during
rendering. Zakharov~et~al.~\cite{Zakharov2022Photo} show that photorealistic
rendering outperforms non-photorealistic alternatives for 6D object detection,
establishing a direct link between rendering fidelity and task performance.
These findings motivate our investigation, yet prior work focuses exclusively
on the gap between simulation and reality. We complement this line of research
by characterizing a \emph{within-simulation} rendering gap: the difference
between full MDL and simplified UsdPreviewSurface materials inside the same
simulator.

\subsection{Material Simplification and Simulation Platforms}

Simplifying complex material models while preserving visual appearance has a
long history in computer graphics, from analytic BRDF
fitting~\cite{Ngan2005Experimental} to data-driven reflectance
compression~\cite{Sztrajman2021Neural}. Meanwhile, simulation platforms for
embodied AI such as Habitat~\cite{Savva2019Habitat},
AI2-THOR~\cite{Kolve2017AI2THOR}, and
ThreeDWorld~\cite{Gan2021ThreeDWorld} typically employ simplified shading models
(e.g., Phong or PBR metallicness) from the outset, avoiding the MDL conversion
problem entirely. In contrast, NVIDIA Isaac Sim couples high-fidelity MDL
materials with RTX ray tracing, making it the dominant platform for robotic
manipulation data generation but also creating a unique conversion challenge
when assets must be exported to lighter-weight renderers. Our work studies this
specific gap between MDL and UsdPreviewSurface within the Isaac Sim ecosystem,
complementing the broader material simplification literature with a focused,
application-driven evaluation.

\subsection{Image Quality Evaluation}

Quantifying visual differences between rendered images requires perceptual
metrics beyond simple pixel error. SSIM~\cite{Wang2004Image} captures
structural distortion, LPIPS~\cite{Zhang2018Unreasonable} measures perceptual
similarity using deep features, and FID~\cite{Heusel2017GANs} quantifies
distributional distance between image populations. While FID and related
distributional metrics are standard for comparing large image populations, they
require sample sizes substantially larger than our $n=24$ pairs to produce
reliable estimates~\cite{Heusel2017GANs}; we therefore focus on paired per-image
metrics. These metrics are routinely
applied to evaluate generative models, but have seen limited use for comparing
rendering configurations within a simulation pipeline. We adopt PSNR, SSIM,
and LPIPS for paired image-level comparison, extending their application to the
material simplification setting.

\subsection{Visual Feature Robustness}

Foundation vision models such as CLIP~\cite{Radford2021Learning} and
DINOv2~\cite{Oquab2023DINOv2} learn representations that generalize across
visual domains and have become standard feature extractors for downstream
tasks. Their robustness to domain shifts makes them natural probes for assessing
whether material simplification alters the semantic content of rendered images.
Unlike pixel-level metrics, cosine similarity in CLIP or DINOv2 embedding space
captures whether high-level visual concepts --- object identity, shape, texture
category --- are preserved despite low-level shading differences. We use both
models to evaluate feature-level preservation as a complement to pixel-level
quality metrics, providing a more complete picture of the simplification
trade-off for AI pipelines.
